\documentclass[11pt]{article}

\usepackage[a4paper,margin=1in]{geometry}
\usepackage{amsmath}
\usepackage{booktabs}
\usepackage{enumitem}
\usepackage[hidelinks]{hyperref}

\title{BayesianAgent: A Partner-Aware OneShot Negotiation Agent}
\author{Kouta Ito}
\date{\today}

\begin{document}
\maketitle

\begin{abstract}
This report describes \texttt{BayesianAgent}, a negotiation agent for the
SCML OneShot setting.  The agent extends a random synchronous OneShot agent
with a lightweight Bayesian-style opponent classifier, first-offer probing,
partner-dependent quantity allocation, and specialized counter-offer handling.
The main design goal is to exploit predictable behavior of
\texttt{GreedyOneShotAgent}-like opponents while retaining robust behavior
against non-greedy partners such as random and equal-distribution agents.
The agent continuously updates a two-class belief over each partner,
distinguishing \texttt{GreedyOneShotAgent} from \texttt{NonGreedy}, and uses
that belief to decide where to send offers, which offers to accept, and when to
send opponent-favorable counter offers.
\end{abstract}

\section{Problem Setting}

In the OneShot track, an agent negotiates independently with multiple
suppliers and consumers at each production step.  Each negotiation outcome is
an offer tuple
\[
  o = (q, t, p),
\]
where \(q\) is the quantity, \(t\) is the delivery step, and \(p\) is the unit
price.  The agent must satisfy its required input and output quantities while
avoiding costly shortage, disposal, or low-profit agreements.  Because
negotiations happen simultaneously, a good policy must reason not only about
single offers but also about combinations of offers from several partners.

The implementation follows the spirit of cautious OneShot agents: it avoids
blindly accepting the set of offers with the smallest absolute quantity
mismatch, and instead treats selling and buying differently.  The current
\texttt{BayesianAgent} combines this cautious quantity matching with a
behavioral model of opponents.  Its most important assumption is that greedy
partners tend to accept offers that are favorable to themselves, while
non-greedy partners often follow distributional or random patterns.

\section{Opponent Model}

\subsection{Two-Class Posterior}

For each partner \(i\), the agent maintains logits for two opponent classes:
\[
  C \in \{\text{GreedyOneShotAgent}, \text{NonGreedy}\}.
\]
The posterior is computed with a softmax:
\[
  P(C=k \mid D_i)
  =
  \frac{\exp(z_{i,k}/\tau)}
       {\sum_{k'} \exp(z_{i,k'}/\tau)},
\]
where \(z_{i,k}\) is the accumulated logit evidence and \(\tau\) is the
temperature parameter \texttt{softmax\_temperature}.  The default temperature is
\(1.0\).  If a partner is vetoed as non-greedy, the posterior is forced to
\[
  P(\text{GreedyOneShotAgent}) = 0,\qquad
  P(\text{NonGreedy}) = 1.
\]

The public classification function \texttt{opponent\_type()} first checks the
number of observations.  If the partner has too few observations, it returns
\texttt{Unknown}.  Otherwise, the most probable class is selected.  A partner is
treated as greedy when \texttt{GreedyOneShotAgent} is the most likely class and
its probability is at least \(0.51\).  Non-greedy classification uses the
general \texttt{classification\_threshold}, whose default value is \(0.60\).

\subsection{Evidence from First Offers}

The classifier is driven mainly by first-offer outcomes.  The agent labels each
price from the opponent's perspective:
\[
\text{opponent score} = 1 - \text{price score for me}.
\]
If the score is at least \(0.80\), the price is labeled \texttt{good}; if it is
at most \(0.20\), it is labeled \texttt{bad}; otherwise it is
\texttt{neutral}.  This convention is important: a \texttt{bad} price means bad
for the opponent and therefore good for the agent.

The main evidence rules are summarized in Table~\ref{tab:evidence}.  A
particularly strong rule is that if the agent sends a first offer with a price
that is bad for the opponent and the partner accepts it, the partner is vetoed
as \texttt{NonGreedy}.  This reflects the idea that a purely greedy opponent
should rarely accept an opponent-unfavorable price.  Conversely, accepting a
good first offer is strong evidence for greedy behavior.

\begin{table}[h]
\centering
\begin{tabular}{lll}
\toprule
Observation & Interpretation & Update \\
\midrule
Bad first offer accepted & opponent accepts unfavorable price & hard non-greedy veto \\
Bad first offer rejected & compatible with greedy behavior & small greedy evidence \\
Good first offer accepted & compatible with greedy behavior & strong greedy evidence \\
Good first offer rejected & suspicious for greedy behavior & small non-greedy evidence \\
Neutral first offer accepted & weak greedy signal & small greedy evidence \\
Neutral first offer rejected & weak non-greedy signal & small non-greedy evidence \\
\bottomrule
\end{tabular}
\caption{Main first-offer evidence used by \texttt{BayesianAgent}.}
\label{tab:evidence}
\end{table}

The agent also records whether non-greedy probes succeed.  These records are
used later to weight offers toward partners that have historically accepted
the agent's first proposals.

\section{First-Proposal Strategy}

\subsection{Exploration Phase}

At the beginning of the game, the agent enters an exploration phase.  This is
controlled by \texttt{exploration\_days}, \texttt{max\_exploration\_days}, and
the ratio of still-unknown partners.  During exploration, the agent distributes
offers over available suppliers and consumers and alternates between prices
that are good and bad for the opponent:
\[
  p =
  \begin{cases}
    \text{worst price for me}, & \text{if probing with an opponent-good price},\\
    \text{best price for me},  & \text{otherwise}.
  \end{cases}
\]
This creates informative accept/reject observations.  A greedy partner should
be more likely to accept an opponent-good price and reject an opponent-bad
price, whereas random or distributional agents may not follow this pattern.

\subsection{Greedy-Aware Allocation}

After exploration, first proposals are generated by
\texttt{\_greedy\_only\_first\_proposals()}.  Partners classified as greedy are
ranked by their greedy posterior.  If there is one greedy partner, the current
implementation sends
\[
  \min(7, \text{needs})
\]
units to that partner at the opponent-good price.  If there are two or more
greedy partners, the agent first computes an \(80\%\) greedy target:
\[
  Q_G = \lceil 0.8 \cdot \text{needs} \rceil.
\]
This target is split between the two highest-probability greedy partners with
an intentional imbalance: an even \(Q_G\) is split with a difference of two
units, while an odd \(Q_G\) is split with a difference of three units.  The
larger quantity is assigned to the partner with the higher greedy probability.
This design avoids sending identical probes to two greedy partners and
emphasizes the more reliable one.

Any remaining target quantity is assigned to non-greedy or unknown partners.
The allocation is weighted by each partner's historical first-offer success
rate:
\[
  w_i = 0.75 + 0.5 \cdot r_i,
\]
where \(r_i\) is the estimated success rate for partner \(i\).  A higher
success rate therefore gives the partner a larger share of the remaining
quantity.  When no partner-specific history is available, the agent falls back
to the global non-greedy success estimate.

\section{Counter-Offer and Acceptance Strategy}

\subsection{Recording Observations}

On every call to \texttt{counter\_all()}, the agent first records incoming
offers and compares them with its latest first proposals.  If a partner replies
to a first proposal with a different offer, the reply is treated as a counter
observation and the corresponding first offer is marked as not accepted.  On
success or failure callbacks, the agent updates the first-offer success
statistics and the posterior evidence.

\subsection{Exact and Near-Exact Quantity Handling}

The agent starts from the response generated by the parent synchronous random
agent, but then rewrites the responses using its own quantity logic.  If the
accepted set already satisfies or exceeds the required quantity, all remaining
negotiations on that side are ended.  Otherwise, the remaining need is
distributed among counter partners.  When only one counter partner is left, the
agent gradually concedes the counter quantity toward the opponent's currently
offered quantity as negotiation time approaches the deadline.

The implementation includes a practical efficiency rule: if the counter offer
would have the same quantity as the incoming offer, the agent accepts the
incoming offer instead of re-offering the same quantity.  This avoids wasting a
negotiation step on an equivalent counter.

\subsection{Seller Greedy-Fill Rule}

The most specialized counter rule is used when the agent is a seller.  If there
is no exact subset of consumer offers whose total equals the needed sales
quantity, and if at least one partner has greedy probability at least
\(\theta_G = 0.51\), the agent selects the most greedy partner as a fill
partner.  All other partners remain eligible for acceptance.  Among these
non-fill partners, the agent selects the subset with the largest total quantity
not exceeding the need:
\[
  S^\star =
  \arg\max_{S}
  \sum_{i \in S} q_i
  \quad
  \text{s.t.}
  \quad
  \sum_{i \in S} q_i \le \text{needs}.
\]
Ties are broken by higher total revenue and then by fewer accepted partners.
The remaining quantity,
\[
  q_{\text{rem}} =
  \text{needs} - \sum_{i \in S^\star} q_i,
\]
is offered to the most greedy partner at the price that is good for the
opponent.  For a seller, this is the worst price for the agent, which is the
price most likely to be accepted by a greedy buyer.  This rule is designed for
cases where accepting many non-greedy offers nearly fills the need, and a
single greedy counter can complete the quantity exactly.

\section{Discussion}

\texttt{BayesianAgent} is built around a simple but useful separation of
concerns.  Classification is based on first-offer behavior, first proposals use
the resulting classification to decide where to place quantity, and counter
responses use subset reasoning to avoid harmful over-acceptance or
under-satisfaction.  The agent is intentionally conservative in its opponent
model: it only distinguishes greedy from non-greedy behavior, avoiding a larger
multi-class model that may overfit sparse observations.

The main strength of this design is adaptivity.  Against greedy agents, the
agent can offer opponent-good prices in exchange for reliable quantity
matching.  Against non-greedy agents, it can spread quantity according to
empirical success rates.  The seller greedy-fill rule further improves
simultaneous offer handling by accepting the best under-need subset and using a
greedy partner to complete the remaining quantity.

The main limitation is that early posterior estimates can be noisy.  At the
beginning of a simulation, many partners have probabilities close to \(0.5\),
so a threshold near \(0.51\) can classify partners aggressively once a small
amount of evidence appears.  This behavior is useful for exploiting greedy
partners quickly, but it may also misclassify random partners.  Future work
could introduce confidence intervals, side-specific thresholds, or separate
models for quantity behavior and price behavior.

\section{Conclusion}

\texttt{BayesianAgent} combines a lightweight Bayesian-style classifier with
partner-aware offer allocation and cautious counter-offer selection.  Its
strategy is designed for the simultaneous nature of OneShot negotiation:
learning which partners are greedy, allocating first proposals according to
that belief, and repairing imperfect offer sets with targeted counters.  The
agent remains simple enough to run inside tournament constraints while
capturing several important behavioral regularities of common OneShot agents.

\begin{thebibliography}{9}
\bibitem{scml}
SCML OneShot and NegMAS documentation and tournament environment.

\bibitem{cautious}
Team Miyajima, \emph{CautiousOneShotAgent} report, ANAC 2024 SCML OneShot
track.
\end{thebibliography}

\end{document}
